% Source: https://artofproblemsolving.com/community/c3h1966987_competition_math_for_middle_school_problem
% Stable identity: aops-topic-1966987-q01
Let $r$ be the number of holes one mole digs per minute. Work done is rate times time times number of moles, $w=rtm$. From the first situation, $4=r\cdot3\cdot5$, so $r=\dfrac{4}{15}$ hole per mole per minute.

For $9$ moles and $6$ holes: $6=\dfrac{4}{15}\cdot t\cdot9=\dfrac{36}{15}t$. To isolate $t$, divide both sides by $\dfrac{36}{15}$, that is, multiply by its reciprocal $\dfrac{15}{36}$:
\[
t=6\cdot\dfrac{15}{36}=\dfrac{90}{36}=\dfrac{5}{2} .
\]
So it takes $\boxed{\dfrac{5}{2}}$ minutes, i.e. $2.5$ minutes. (Check: $9$ moles for $2.5$ minutes is $22.5$ mole-minutes, and $6$ holes need $6\cdot\dfrac{15}{4}=22.5$ mole-minutes.)
